\documentclass[11pt]{ltjsarticle}
\usepackage{luatexja-fontspec}
\usepackage{amsmath,amssymb}
\usepackage[margin=22mm]{geometry}
\usepackage{graphicx}
\usepackage{longtable,booktabs,array}
\usepackage{calc}
\usepackage{xcolor}
\usepackage{hyperref}
\hypersetup{colorlinks=true,linkcolor=blue,urlcolor=blue,citecolor=blue}
\makeatletter
\def\maxwidth{\ifdim\Gin@nat@width>\linewidth\linewidth\else\Gin@nat@width\fi}
\makeatother
\setkeys{Gin}{width=\maxwidth,keepaspectratio}
\setcounter{secnumdepth}{0}
\setlength{\parindent}{0pt}
\setlength{\parskip}{0.5em}
\providecommand{\tightlist}{\setlength{\itemsep}{0pt}\setlength{\parskip}{0pt}}
\providecommand{\pandocbounded}[1]{#1}
\providecommand{\real}[1]{#1}
\newcounter{none}
\begin{document}
\section{補遺：複素構造は措定でなく帰結である}\label{ux88dcux907aux8907ux7d20ux69cbux9020ux306fux63aaux5b9aux3067ux306aux304fux5e30ux7d50ux3067ux3042ux308b}

\subsection{------二つの実成分から複素平面が立ち上がること}\label{ux4e8cux3064ux306eux5b9fux6210ux5206ux304bux3089ux8907ux7d20ux5e73ux9762ux304cux7acbux3061ux4e0aux304cux308bux3053ux3068}

\textbf{著者：} 木原 範昭（ORCID: 0009-0004-6753-4020） \textbf{日付：}
2026年6月18日 \textbf{DOI（本版）：} 10.5281/zenodo.20742985 ／
\textbf{Concept DOI（全版）：} 10.5281/zenodo.20742984 ／
\textbf{Zenodo:} https://zenodo.org/records/20742985
\textbf{関連論文：}「高い対称性の極限におけるスケールの対数写像とスケール・量子化の分離について」§4・§7.4（Concept
DOI: 10.5281/zenodo.20740841）

\begin{center}\rule{0.5\linewidth}{0.5pt}\end{center}

\subsection{趣旨}\label{ux8da3ux65e8}

本補遺は、関連論文 §4 および §7.4
の注で数行に圧縮した一点------\textbf{複素数 i
を公理に置かず、二つの実成分と一つの回転から、複素構造が帰結として現れる}こと------を、独立に展開する。親論文に新しい主張を加えるものではなく、誤読されやすい「i
がどこから来るか」を明示することを目的とする。

通常、量子論は複素ヒルベルト空間を基本枠組みとして与え、虚数単位 i
を最初に措定する。本系では i
を公理に置かない。公理に置かれるのは、二つの独立な実保存量------二乗和（差異の束ね、g）と積（在る/無いの縁、ω）------と、それらを結ぶ回転だけである。複素構造（複素平面・i・複素内積）は、この実数二者性の帰結として立ち上がる。本補遺はその立ち上がりを記述する。

本補遺は厳密な証明ではなく、親論文の構造を明示する観察である。物理量への同定は行わない。本補遺は「量子論が複素数で記述される理由を解明した」と主張するものではなく、本系の最小公理から複素構造が形式的に現れる経路を示すに留まる。

\begin{center}\rule{0.5\linewidth}{0.5pt}\end{center}

\subsection{1. 出発点：二つの実成分、i
は無い}\label{ux51faux767aux70b9ux4e8cux3064ux306eux5b9fux6210ux5206i-ux306fux7121ux3044}

本系の公理は実数で閉じる。複素数は導入しない。基本となるのは、二つの独立な実保存量である（親論文
公理4・公理5）。

二つの実成分を ν\textsubscript{1}, ν\textsubscript{2}
と書く。これらは無名であり、どちらを第一・第二と呼ぶかは任意である（親論文の無名性）。両者は対等な二つの実数の軸をなす。

虚数単位 i は、この段階では存在しない。以下で示すのは、ν\textsubscript{1}, ν\textsubscript{2}
という二つの実成分の上に、二乗和と積という二つの保存量が立ち、それらが複素数の実部・虚部と同じ代数構造をなす、ということである。

\begin{center}\rule{0.5\linewidth}{0.5pt}\end{center}

\subsection{2.
二つの保存量：二乗和（g）と積（ω）}\label{ux4e8cux3064ux306eux4fddux5b58ux91cfux4e8cux4e57ux548cgux3068ux7a4dux3c9}

二つの実成分 ν\textsubscript{1}, ν\textsubscript{2} に対し、二つの独立な保存量を考える。

\textbf{二乗和（差異の束ね、g）：}

\[g = \nu_1^2 + \nu_2^2.\]

これは二者を対等に束ねる量であり、ラベルの読み替え（ν\textsubscript{1}↔ν\textsubscript{2}）に不変である。SO(2)------二乗和を保つ回転------の不変量である（親論文
公理5）。

\textbf{積（在る/無いの縁、ω）：} 揺らぎの積

\[\omega = \delta_1\delta_2 = \tfrac{1}{2} \quad(\text{面積・シンプレクティック形式；親論文 公理4}).\]

これは面積（二軸が張る量）であり、SO(1,1)------双曲的な、積を保つ変換------に関係する、回転で保たれる不変量である（親論文
公理4）。なお成分そのものの積 ν\textsubscript{1}ν\textsubscript{2} は回転 θ
に依存する別の双線形量であり、ω（δ\textsubscript{1}δ\textsubscript{2}）そのものではない（§3）。

親論文の注（公理4 と公理5 の相互導出不能）の通り、g と ω
は互いに導出できない。「在るか」（積・ω）と「違うか」（二乗和・g）は、無名性の下で独立に立つ二つの根源区別だからである。

\begin{center}\rule{0.5\linewidth}{0.5pt}\end{center}

\subsection{3. 回転：cos と sin
の導入}\label{ux56deux8ee2cos-ux3068-sin-ux306eux5c0eux5165}

二つの実成分を、動径 ν と角 θ
で表す。これは二次元の二成分を極形式で書くだけであり、新しい仮定ではない：

\[\nu_1 = \nu\cos\theta, \qquad \nu_2 = \nu\sin\theta.\]

cos と sin は、二乗和を一定に保つ回転（SO(2)）のパラメータである。ここで
cos²θ+sin²θ=1 により、

\[g = \nu_1^2 + \nu_2^2 = \nu^2(\cos^2\theta + \sin^2\theta) = \nu^2\]

となり、二乗和 g は動径の二乗 ν² に等しく、回転 θ
に依らず一定である（動径のみで決まる）。一方、個々の成分の積は

\[\nu_1\nu_2 = \nu^2\cos\theta\sin\theta = \tfrac{1}{2}\nu^2\sin 2\theta\]

となり、θ とともに変化する。これは成分レベルの双線形量であって、ω
そのものではない。複素構造の虚部に対応する面積
ω（δ\textsubscript{1}δ\textsubscript{2}=1/2、§2）は、回転で保たれるシンプレクティックな不変量である。

\begin{center}\rule{0.5\linewidth}{0.5pt}\end{center}

\subsection{4.
複素構造の立ち上がり}\label{ux8907ux7d20ux69cbux9020ux306eux7acbux3061ux4e0aux304cux308a}

ここで、二つの実成分 (ν\textsubscript{1}, ν\textsubscript{2}) を一つの対象として束ねる。束ねた対象を z
と書き、その二成分を「実部」「虚部」と呼ぶ：

\[z = \nu_1 + i\,\nu_2 = \nu(\cos\theta + i\sin\theta) = \nu\,e^{i\theta}.\]

ここで現れた i は、措定されたものではない。\textbf{i
は、二つの実軸を結ぶ 90 度回転 J の符牒}である（§3 の
cos↔sin）。二者性（独立な2軸）は i
が作用する「平面」を与え、その平面上の 90 度回転 J が i
に対応する。\textbf{J²=−1}------90 度回転を 2 回行えば 180 度回転＝−1
倍------であることが、i²=−1 という規則になる。i
は外から導入した新しい数ではなく、二軸の上の回転を代数的に書いた記号である。（独立な
2 軸だけでは実平面 $\mathbb{R}$² にすぎず、回転 J を与えて初めて複素平面 $\mathbb{C}$
になる。）

この束ねの下で、三つの対応が立つ：

\begin{itemize}
\tightlist
\item
  \textbf{二乗和 g = ν\textsubscript{1}²+ν\textsubscript{2}² = \textbar z\textbar²}
  ------複素数の絶対値の二乗（実部）。
\item
  \textbf{面積 ω = δ\textsubscript{1}δ\textsubscript{2}}
  ------複素数の虚部側に属する量（シンプレクティック形式）。
\item
  \textbf{回転 cos↔sin} ------複素構造 J（90度回転、i
  を掛けることに対応）。
\end{itemize}

すなわち、二つの実成分・二乗和・積・回転という、すべて実数で書かれた構造が、複素数
z=νe\^{}\{iθ\} とその代数（絶対値・偏角・i
倍）と一対一に対応する。複素平面は、実数二者性の上に立ち上がる像である（図1）。

\textbf{図1（fig3\_complex\_from\_two\_real.svg）：}
二つの実成分から複素構造が立ち上がる。実軸 ν\textsubscript{1}・ν\textsubscript{2} の上に点 z
が乗り、動径 ν=√g は回転に不変（破線の円）、個々の成分は角 θ
とともに変化する。この対 (ν\textsubscript{1}, ν\textsubscript{2}) が複素数 z=νe\^{}\{iθ\} に対応し、i
は二軸を結ぶ 90 度回転 J（J²=−1）の記号として現れる。

\textbf{複素内積との対応。} 二つの保存量を一つの複素量に束ねると、

\[\langle\,\cdot\,,\cdot\,\rangle = g + i\,\omega\]

という形が得られる。実部 g（二乗和・距離・計量）と虚部
ω（積・面積・シンプレクティック形式）が、複素内積の実部・虚部に対応する。親論文
§7 の g/ω
分離は、この複素内積の実部と虚部が、対数写像の下で別々の運命をたどる（g
は平坦化、ω は不変）ことを述べたものである。

\begin{center}\rule{0.5\linewidth}{0.5pt}\end{center}

\subsection{5. 構造の要約}\label{ux69cbux9020ux306eux8981ux7d04}

{\def\LTcaptype{none} % do not increment counter
\begin{longtable}[]{@{}ll@{}}
\toprule\noalign{}
実数で書いた構造 & 複素表現での対応 \\
\midrule\noalign{}
\endhead
\bottomrule\noalign{}
\endlastfoot
二つの実成分 ν\textsubscript{1}, ν\textsubscript{2} & 複素数 z の実部・虚部 \\
二軸を結ぶ 90 度回転 J & i（J²=−1 すなわち i²=−1） \\
二乗和 g = ν\textsubscript{1}²+ν\textsubscript{2}² & \\
面積 ω = δ\textsubscript{1}δ\textsubscript{2} & 虚部側（シンプレクティック形式） \\
回転 cos↔sin（SO(2)） & 複素構造 J（i 倍、偏角 θ） \\
二つの保存量の束ね & 複素内積 $\langle$·,·$\rangle$ = g + iω \\
\end{longtable}
}

複素数は公理に置かれていない。公理に置かれるのは、二つの独立な実保存量（二乗和
g・積 ω）と回転だけである。i は、二軸を結ぶ 90
度回転（J）を表す記号として、束ねの段階で現れる。複素平面・複素内積は、この実数二者性（親論文の「無名性の限界＝2」が与える二者性）と回転の帰結である。

なぜ底が 2 か、なぜ二者か------それは親論文
公理2（無名性の限界＝2）による。系の中で差異が立つ最小かつ上限が二であること、これが「二つの実成分」の根拠であり、複素構造が二次元（一つの実部・一つの虚部）である根拠でもある。複素数が二次元であることは、無名性の限界が
2 であることの反映である。

\begin{center}\rule{0.5\linewidth}{0.5pt}\end{center}

\subsection{6.
限界（読み解きの注意）}\label{ux9650ux754cux8aadux307fux89e3ux304dux306eux6ce8ux610f}

本補遺は親論文 §4・§7.4
の構造の明示であり、新しい主張を加えない。以下を明記する。

\begin{enumerate}
\def\labelenumi{\arabic{enumi}.}
\item
  \textbf{本補遺は「量子論が複素数で記述される理由」を解明したと主張しない。}
  本系の最小公理から複素構造が形式的に現れる一つの経路を示すに留まる。物理理論の複素性との同型は別途の課題である。
\item
  \textbf{i は二軸を結ぶ 90 度回転 J の記号として現れる（J²=−1）。}
  これは「i
  に物理的実在を与えた」という主張ではなく、二実成分の上の回転を一つの複素量に束ねる代数的記法の帰結である。独立な
  2 軸だけでは実平面 $\mathbb{R}$² にすぎず、回転 J を加えて初めて複素平面 $\mathbb{C}$
  になる。実数二者性と回転を複素記法で書くか、二実成分＋回転で書くかは表現の選択であり、本系の内容は実数で閉じている。
\item
  \textbf{二乗和 g と積 ω の相互導出不能は前提である。}
  親論文の注の通り、g と ω
  は独立な二つの根源区別であり、一方から他方を導いていない。複素構造は、この二つが独立に立つことの上に現れる。
\item
  \textbf{物理的同定を行わない。} 二つの実成分 ν\textsubscript{1}, ν\textsubscript{2}
  の物理的名称は無名のラベルであり、本補遺は特定の物理量への同定を行わない。
\item
  \textbf{二者性の根拠は親論文 公理2 に依存する。}
  「なぜ二つの実成分か（なぜ複素数が二次元か）」は、無名性の限界＝2（親論文
  公理2）に帰せられる。この公理自体の導出は親論文・本補遺の射程外である。
\end{enumerate}

\begin{center}\rule{0.5\linewidth}{0.5pt}\end{center}

\subsection{参照}\label{ux53c2ux7167}

\begin{itemize}
\tightlist
\item
  Kihara,
  N.「高い対称性の極限におけるスケールの対数写像とスケール・量子化の分離について------最小公理系からの限定的検証報告（改訂版）」§4・§7.4（複素構造は実数二者性の帰結）、公理4・公理5。Concept
  DOI: 10.5281/zenodo.20740841。
\end{itemize}
\end{document}
